% Abstract + Section I. Budgets: abstract 0.4 col (~180 words), intro 0.8 col (PAPER_PLAN.md).
\begin{abstract}
FPGA accelerator papers often report performance from models or simulation,
leaving the gap to real hardware unquantified. We present an INT8 CNN
accelerator for the AMD Kria KV260 whose performance is exactly predictable.
An $8{\times}8$ output-stationary array of DSP48E2 slices runs LeNet-5 and a
CIFAR-10 network entirely on chip at 250\,MHz, using 80 DSPs and 16.5k LUTs.
Its integer requantization reproduces the floating-point reference on every
non-saturating input of every channel. An analytical cycle model, fixed
before RTL simulation, predicts every layer's cycle count with zero error on
20,000 test images and 286 random layer configurations, in RTL simulation
and on the board. Measured on the KV260, accelerator compute takes 65.7\,$\mu$s
(LeNet-5) and 417.0\,$\mu$s (CIFAR-10) with a single cycle count over 3.65M
inferences, and 0.065 and 0.161\,mJ per image of SOM-rail energy. End to
end through a Python host it is 2.1--3.1$\times$ faster than the best
on-board Cortex-A53 baseline; against the AMD DPU it is 1.5$\times$ faster on
LeNet-5 and 1.6$\times$ slower on CIFAR-10. All results are regenerated by
script from a tagged commit.
\end{abstract}
% src: tab_t1_impl (80 DSP, 16,538 LUT @250); tab_layer_spread / tab_shapes (board, cycle-exact);
% tab_soak (3,651,181 jobs, 1 distinct TOTAL_CYC); tab_a5_cpu (65.7 / 417.0 us; speedups 3.10x / 2.11x;
% DPU/accel e2e 1.54x / 0.61x); tab_b1_power (E_sys 0.065 / 0.161 mJ).

\section{Introduction}\label{sec:intro}
FPGA CNN accelerators are usually evaluated with an analytical model or a
simulator, and the design that reaches the board is then characterized by
a few throughput figures. How far the model is from the hardware, layer by
layer, is rarely reported. That gap matters for two reasons. Embedded and
real-time systems need latency they can budget for, not latency that is
usually close. And a performance model is only useful for design-space
exploration if it can be trusted without building every design.

This paper takes the opposite route to the usual one. We wrote the cycle
model first, derived its two constants from the RTL's stage latencies and
committed them before the first simulation of the core, and then held the
RTL and the board to that model cycle for cycle. The design is deliberately
small: one $8{\times}8$ INT8 array on the AMD Kria KV260, all weights in BRAM,
and two small networks. Our claim is exactness and determinism, not peak
throughput, and we compare it honestly with the CPU and the vendor DPU on
the same board. Our contributions are:
\begin{itemize}
\item An $8{\times}8$ output-stationary INT8 accelerator on the KV260 whose
integer requantization is checked bit-exact against the floating-point
reference on every non-saturating input of every channel
(Secs.~\ref{sec:arch}, \ref{sec:verif}).
\item A closed-form cycle model, with constants fixed before RTL simulation,
that predicts every layer's cycles exactly on 20,000 test images, 286 random
layer jobs and the board (Secs.~\ref{sec:model}, \ref{sec:results}).
\item A board evaluation on the same KV260: deterministic compute latency,
SOM-rail energy, a 7-point clock sweep, and a comparison with the on-board
Cortex-A53 and the AMD DPU (Sec.~\ref{sec:results}).
\item A reproducible artifact: every reported number is generated by script
from a tagged commit, with hashed logs and bitstreams.
\end{itemize}
